\documentclass[aspectratio=169,UTF8,fontset=fandol]{ctexbeamer}
\usepackage{amsmath}
% Preamble fragment. The containing document owns its fonts and page geometry.
\usepackage{tikz}
\usetikzlibrary{arrows.meta,positioning,calc,shapes.geometric,fit,backgrounds}
\definecolor{sqInk}{HTML}{243247}
\definecolor{sqBlue}{HTML}{245E95}
\definecolor{sqTeal}{HTML}{167D7D}
\definecolor{sqAmber}{HTML}{9A5E16}
\tikzset{
  sq base/.style={font=\normalsize,text=sqInk,align=center,inner sep=5pt},
  sq process/.style={sq base,rounded corners=2pt,draw=sqBlue,
    fill=sqBlue!5,line width=0.8pt,text width=2.45cm,minimum height=1.0cm},
  sq result/.style={sq process,draw=sqTeal,fill=sqTeal!6},
  sq pending/.style={sq process,draw=sqAmber,fill=sqAmber!7},
  sq decision/.style={sq base,font=\small,diamond,aspect=1.8,draw=sqBlue,
    fill=sqBlue!5,line width=0.8pt,text width=2.15cm,inner sep=2pt},
  sq line/.style={draw=sqInk,line width=0.85pt},
  sq arrow/.style={sq line,-{Stealth[length=2mm,width=1.4mm]}},
  sq label/.style={font=\small,text=sqInk,inner sep=1.5pt}
}

\setbeamertemplate{navigation symbols}{}
\setbeamercolor{frametitle}{fg=sqBlue}
\setbeamertemplate{footline}[frame number]
\begin{document}
\begin{frame}{从观测到候选：求解条件不等于完成证明}
  \centering
  % Body fragment; requires diagram-styles.tex in the preamble.
\begin{tikzpicture}[node distance=0.72cm]
  \node[sq process,text width=2.1cm] (observations)
    {示范数据\\$(x_i,u_i)$};
  \node[sq process,text width=2.1cm,right=of observations] (features)
    {构造特征\\$\Psi(x,u)$};
  \node[sq process,text width=2.1cm,right=of features] (conditions)
    {辨识条件\\$G\theta=b$};
  \node[sq result,text width=2.1cm,right=of conditions] (candidates)
    {候选代价\\$\theta\in\Theta$};
  \draw[sq arrow] (observations.east) -- node[sq label,above] {估计} (features.west);
  \draw[sq arrow] (features.east) -- node[sq label,above] {堆叠} (conditions.west);
  \draw[sq arrow] (conditions.east) -- node[sq label,above] {求解} (candidates.west);
\end{tikzpicture}

  \par\vspace{0.7cm}
  特征模型与最优性条件成立时，构造候选集合。\\[0.2cm]
  唯一性、可行性和全局最优性仍需另外验证。
\end{frame}
\begin{frame}{有边界的修订：验证通过或明确报告未通过}
  \centering
  % Body fragment; the two exits distinguish acceptance from exhausted budget.
\begin{tikzpicture}
  \node[sq process] (candidate) at (-4.2,1.4) {构建候选\\保存当前版本};
  \node[sq decision] (validation) at (0,1.4) {验证通过？\\$\varepsilon\leq\tau$};
  \node[sq result] (accepted) at (4.2,1.4) {交付当前版本\\附验证记录};
  \node[sq decision] (budget) at (0,-1.1) {修订预算\\仍有剩余？};
  \node[sq process] (revision) at (-4.2,-1.1) {修订候选\\消耗一次预算};
  \node[sq pending] (unresolved) at (4.2,-1.1) {报告未通过\\保留现有证据};
  \draw[sq arrow] (candidate.east) -- (validation.west);
  \draw[sq arrow] (validation.east) -- node[sq label,above] {是} (accepted.west);
  \draw[sq arrow] (validation.south) -- node[sq label,right] {否} (budget.north);
  \draw[sq arrow] (budget.west) -- node[sq label,above] {是} (revision.east);
  \draw[sq arrow] (budget.east) -- node[sq label,above] {否} (unresolved.west);
  \draw[sq arrow] (revision.north) -- (candidate.south);
\end{tikzpicture}

  \par\vspace{0.35cm}
  未通过验证时，仅在预算允许的范围内继续修订。
\end{frame}
\begin{frame}{负反馈：控制、扰动与测量分别流向哪里}
  \centering
  % Body fragment; scalar negative-feedback schematic, not a stability claim.
\begin{tikzpicture}
  \node[sq base] (reference) at (-5.7,0) {$r_t$};
  \node[circle,draw=sqInk,line width=0.85pt,minimum size=0.6cm,inner sep=0pt]
    (sum) at (-4.4,0) {};
  \node[sq process,text width=1.8cm] (controller) at (-1.9,0) {控制器\\$\pi$};
  \node[sq process,text width=1.8cm] (plant) at (1.8,0) {被控系统\\$f$};
  \coordinate (tap) at (4.1,0);
  \node[sq base] (measurement) at (5.1,0) {$y_t$};
  \node[sq base] (disturbance) at (1.8,1.8) {扰动 $d_t$};
  \draw[sq arrow] (reference.east) -- (sum.west);
  \node[font=\small] at ($(sum.west)+(-0.16,0.22)$) {$+$};
  \node[font=\small] at ($(sum.south)+(0.22,-0.12)$) {$-$};
  \draw[sq arrow] (sum.east) -- node[sq label,above] {$e_t$} (controller.west);
  \draw[sq arrow] (controller.east) -- node[sq label,above] {$u_t$} (plant.west);
  \draw[sq arrow] (plant.east) -- (measurement.west);
  \draw[sq arrow] (disturbance.south) -- (plant.north);
  \fill[sqInk] (tap) circle (1.5pt);
  \draw[sq arrow] (tap) -- ++(0,-1.65) -| (sum.south);
  \node[sq base,font=\small] at (0,-1.95) {测量反馈 $y_t$，误差 $e_t=r_t-y_t$};
\end{tikzpicture}

  \par\vspace{0.35cm}
  这是信号关系示意，不代表系统已被证明稳定。
\end{frame}
\begin{frame}{同一候选的两类检查，支持不同判断}
  \centering
  % Body fragment; shared candidate, distinct checks, no assumed pass verdict.
\begin{tikzpicture}
  \node[sq process] (candidate) at (-4.7,0) {同一候选\\$\theta$};
  \coordinate (fork) at (-2.7,0);
  \node[sq process] (fit) at (0,1.2) {拟合检查\\$\|G\theta-b\|$};
  \node[sq process] (control) at (0,-1.2) {重解控制\\$\hat\pi_\theta$};
  \coordinate (join) at (2.7,0);
  \node[sq result] (evidence) at (4.7,0) {对照检查证据\\记录适用范围};
  \draw[sq line] (candidate.east) -- (fork);
  \draw[sq arrow] (fork) |- (fit.west);
  \draw[sq arrow] (fork) |- (control.west);
  \draw[sq line] (fit.east) -| (join);
  \draw[sq line] (control.east) -| (join);
  \draw[sq arrow] (join) -- (evidence.west);
\end{tikzpicture}

  \par\vspace{0.35cm}
  拟合误差小，不自动意味着重解控制与专家一致。\\[0.15cm]
  两条检查路径不表示实际执行一定并发。
\end{frame}
\end{document}
